% TT09031C
% Rückenmarksverletzung, Verdacht
% 14.0
% 2013-11-30
:ref:`m-ws-trauma-verdacht`
Taktik 
Immobilisation und Transport in eine geeignete Einrichtung
zur Abklärung}


Vgl. \prettyref{m:einschaetzung-indikation-ws-immobilisation}
\begin{itemize}
  -   \EE{\Abk{HWS}-Immobilisation} (HWS-Schiene)
  -   Bergung mit Schaufeltrage, Spineboard oder
  KED-System
  -   **Lagerung** und Schienung mittels Vakuummatratze, Schaufeltrage mit
  Fixationsmöglichkeit oder Spineboard.\footnote{Welches Hilfsmittel
  eingesetzt wird ist abhängig vom Patientenzustand,
  dem Training des Teams, internen Vorschriften und
  vorhandenem Material.}
  -   **Monitoring**: Achte besonders auf Störungen
  von Atmung und Bewusstsein, Schockentwicklung und
  beginnende neurologische Ausfälle.
%   -   **Unfallmechanismus** erheben und dokumentieren
  -   **Transportentscheidung**: Je nach
  lokalen Gegebenheiten
\end{itemize}